\documentclass{article}%
\usepackage[T1]{fontenc}%
\usepackage[utf8]{inputenc}%
\usepackage{lmodern}%
\usepackage{textcomp}%
\usepackage{lastpage}%
\usepackage{geometry}%
\geometry{margin=2cm}%
\usepackage{expex}%
\usepackage[T1]{fontenc}%
\usepackage[utf8]{inputenc}%
\usepackage[spanish]{babel}%
%
\title{Análisis Lingüístico: Gramatica Normativa Kaqchikel}%
\author{Generado por el TFM de Elías Carrasco}%
%
\begin{document}%
\normalsize%
\maketitle%
\section{Page 170}%
\label{sec:Page170}%
plural) del objeto y sujeto. La palabra principal incluye primeramente el tiempo/aspecto %
o modo (excepto en el perfectivo que es la última), seguido por el juego absolutivo, %
juego ergativo, la raíz o base en calidad de obligatorios y sufijos de categoría en casos %
específicos (si el verbo es derivado). Como elementos opcionales incluye el movimiento, %
su posición es siempre antes del juego ergativo. Ejemplos: %
\subsection{Marca de tiempo/aspecto/modo:}%
\label{subsec:Marcadetiempo/aspecto/modo}%

%
\par\medskip%
\ex
\textit{Todavía fui a ver al niño.} \\
\begingl
\gla Xb'entz'eta' na kan ri ak'wal. //
\glb Xb'entz'et a' //
\glc \textbf{SVT} \textbf{SVT} //
\endgl
\xe
%
\medskip%
\par\medskip%
\ex
\textit{} \\
\begingl
\gla X-0-b'e-n-tz'et-a' //
\glb X-0-b'e-n-tz'et-a' //
\glc \textbf{COM-}\textbf{B3s-}\textbf{MOV-}\textbf{Als-}\textbf{ver-}\textbf{SC} //
\glft \textbf{Sintaxis:} [na]\textsubscript{PAR} [kan]\textsubscript{DIR} [ri]\textsubscript{DET} [ak'wal]\textsubscript{niño} //
\endgl
\xe
%
\medskip%
\subsection{Marca el objeto en el verbo:}%
\label{subsec:Marcaelobjetoenelverbo}%

%
\par\medskip%
\ex
\textit{...vio muy hermosas a las niñas...} \\
\begingl
\gla ...kan jeb'ël ok xerutz'ët ri ch'uta'q xtani' rija'... //
\glb kan //
\glc \textbf{SVT} //
\glft \textbf{Sintaxis:} [jeb'ël]\textsubscript{ok} [xerutz'ët]\textsubscript{ri} [ch'uta'q]\textsubscript{xtani'} [rija'...]\textsubscript{SVT} //
\endgl
\xe
%
\medskip%
\par\medskip%
\ex
\textit{} \\
\begingl
\gla ...kan jeb'ël ok x-e-ru-tz'ët ri ch'uta'q xtan-i' rija'... //
\glb kan jeb'ël ok x-e-ru-tz'ët ri ch'uta'q xtan-i' rija'... //
\glc \textbf{PAR} \textbf{hermosa} \textbf{DIR} \textbf{COM-}\textbf{B3p-}\textbf{A3s-}\textbf{ver} \textbf{DET} \textbf{DIM} \textbf{Niña-} \textbf{PL} //
\endgl
\xe
%
\medskip%
\subsection{Marca el sujeto en el verbo:}%
\label{subsec:Marcaelsujetoenelverbo}%

%
\par\medskip%
\ex
\textit{Los dos hombres compraron un corte.} \\
\begingl
\gla Xkilöq' jun uqaj ri ka'i' achi'a'. //
\glb Xkilöq' //
\glc \textbf{SVT} //
\glft \textbf{Sintaxis:} [jun uqaj]\textsubscript{ri} [ka'i' achi'a'.]\textsubscript{IND} //
\endgl
\xe
%
\medskip%
\par\medskip%
\ex
\textit{} \\
\begingl
\gla X-0-ki-löq //
\glb X-0-ki-löq //
\glc \textbf{COM-}\textbf{B3s-}\textbf{A3p-}\textbf{comprar} //
\glft \textbf{Sintaxis:} [jun uqaj]\textsubscript{IND} [ri]\textsubscript{DET} [ka'i']\textsubscript{dos} [achi'-a'.]\textsubscript{hombre-PL} //
\endgl
\xe
%
\medskip%
\par\medskip%
\ex
\textit{La señorita regaló el tomate.} \\
\begingl
\gla Xusipaj ri xkoya' ri xtän. //
\glb Xusipaj //
\glc \textbf{SVT} //
\glft \textbf{Sintaxis:} [ri]\textsubscript{DET} [xkoya']\textsubscript{ri} [xtän.]\textsubscript{DET} //
\endgl
\xe
%
\medskip%
\par\medskip%
\ex
\textit{} \\
\begingl
\gla X-0-u-sip-a-j //
\glb X-0-u-sip-a-j //
\glc \textbf{COM-}\textbf{B3s-}\textbf{A3s-}\textbf{regalar-}\textbf{BV-}\textbf{SC} //
\glft \textbf{Sintaxis:} [ri]\textsubscript{DET} [xkoya']\textsubscript{tomate} [ri]\textsubscript{DET} [xtän.]\textsubscript{señorita} //
\endgl
\xe
%
\medskip%
\subsection{Marca movimiento en el verbo:}%
\label{subsec:Marcamovimientoenelverbo}%

%
\par\medskip%
\ex
\textit{...nuestro papá va a trabajar, vamos a ayudarlo....} \\
\begingl
\gla ...nb'esamäj qatata', nb'eqato'... //
\glb nb'esamäj //
\glc \textbf{SVT} //
\glft \textbf{Sintaxis:} [nb'eqato']\textsubscript{SVT} //
\endgl
\xe
%
\medskip

%
\end{document}